%=================================%
\section{Introduction}
\label{sec:background}
%=================================%


Gravitational coupling between a gaseous disc and an embedded massive perturber --- a planet or a satellite --- is known to be one of the natural drivers of the protoplanetary disc evolution and a key factor determining the architectures of exoplanetary systems \citep{Paardekooper20232023}. Since the early pioneering studies of \citet{GoldreichTremaine_1979,GoldreichTremaine_1980} and \citet{Lin1979} it was known that the exchange of angular momentum between a planet and its nascent disc can lead to planet migration and may truncate the disc at certain locations leading to gap opening. The latter effect may ultimately be responsible for the origin of various substructures that have been observed in protoplanetary discs through e.g. thermal dust emission \citep{AndrewsEtAl_2018,Andrews2020}.

It is also well known observationally that orbits of many exoplanets have appreciable eccentricities \citep{Santos2007,Xie2016}. While the mechanisms and timing of the eccentricity production are still debated \citep{Davies2014}, it is certainly plausible that planets could have a significant eccentricity while still embedded in protoplanetary discs, e.g. due to planet-planet scattering\footnote{More massive planets capable of opening clean gaps may develop non-zero eccentricity also due to gravitational planet-disc coupling \citep{Goldreich2003,Ragusa2018}.}. This naturally raises a possibility of gravitational coupling between the disc and an embedded eccentric planet, a topic on which a number of results have been obtained since the early days in the limit of small planetary eccentricity \citep{GoldreichTremaine_1980,Ward_1986,Artymowicz_1993}. However, the approximations often made in these studies and the differences in their predictions \citep{PapaloizouLarwood_2000,CresswellNelson_2006,IdaEtAl_2020} warrant a re-evaluation of this problem using a first-principles calculation. This is what we do in this work, focusing on the limit of low mass planets to make our calculations rigorous.

Planet-disc coupling is a two-stage process mediated by the density waves driven in the disc by the planet. In the first stage the non-axisymmetric planetary potential excites a density wave, injecting angular momentum into the wave via the {\it excitation} torque acting on the perturbed surface density in the disc. From the global angular momentum conservation, an equal and opposite torque must act on the planet, driving the evolution of its orbital elements --- semi-major axis and eccentricity (in the planar case). In the second stage, the wave damps as it propagates through the disc, either linearly \citep{Takeuchi1996,MirandaRafikov_2020} or non-linearly \citep{GoodmanRafikov_2001,Rafikov2002a}, ultimately depositing its angular momentum into the disc fluid. It is this latter process that causes the disc to evolve, leading to gap opening and the emergence of various observable substructures. 

In the present work (in line with many existing studies) we will be interested only in the former, wave excitation, process, which is important to understand even when it is considered in isolation. First, in most cases the excitation torque fully determines the orbital evolution of planets, which is the main motivation for this study. Second, the excitation torque should be quite insensitive to the details of the subsequent wave damping unless the wave dissipation is very fast (a possibility which we do not consider in this study). For brevity, in the following we will refer to the excitation torque simply as `torque'. There are multiple ways\footnote{Some studies, e.g. \citet{IdaEtAl_2020}, have synthesized orbital evolution prescriptions using multiple approaches listed below.} in which planetary torque and the evolution of planetary orbital elements can be explored. 

First, they can be extracted directly from numerical simulations of eccentric planet-disc interaction, either in 2D or 3D geometry \citep{CresswellNelson_2006,CresswellNelson_2008,Cresswell2007} by following the obital evolution of simulated planets gravitationally interacting with the disc. 

Second, there have been attempts to compute the planetary torque by considering the force acting on the planet as a variant of the dynamical friction against the gaseous background of the underlying disc \citep{MutoEtAl_2011,Desj2022,BuehlerEtAl_2024,ONeill2024}. Unfortunately, to make the calculation tractable these studies have to assume an idealized unperturbed background state which differs from the sheared Keplerian flow in a real disc. 

Third, many studies of the planetary orbital evolution employ a Fourier-space based approach rooted in the analytical methodology of the torque calculation proposed in the seminal works of \citet{GoldreichTremaine_1979,GoldreichTremaine_1980} and \citet{Lin1979}. The idea lies in linearization and (azimuthal) Fourier decomposition of the fluid equations governing planet-disc coupling, with the resultant equations possessing a number of singularities at the radial locations corresponding to the Lindblad and corotation resonances (see Section \ref{subsection:resonances}). A subsequent simplification of these equations in the vicinity of these resonances in the WKB limit shows that (1) the torque is excited only at these resonances and (2) its amplitude can be related in a very simple way to the planetary potential (and its radial derivative) evaluated at the resonance, {\it without actually solving} the original equations for fluid perturbations. Subsequent studies introduced various modifications to this prescription, which incorporated the leading order background gradients in the disc \citep{Ward_1986,Ward_1988}, abandoned the WKB assumption and reintroduced the effect of azimuthal pressure gradients for large azimuthal wavenumbers \citep{Artymowicz_1993, Artymowicz1993b, Artymowicz_1994, Ward_1997}. The latter refinements modified the torque formula to sample the perturbing potential over a finite radial interval around the resonance (rather than at its exact location) and allowed it to account for a \textit{torque cut-off} \citep[first described in][]{GoldreichTremaine_1980} which quenches the torque contributions from the highest order modes. Unfortunately, results obtained with this approach are only approximate even in the linear limit as a result of local, near-resonance expansion of the governing equations. Furthermore, they do not easily yield the radial distribution of the planetary torque across the disc \citep[only its integrated value, cf.] []{RafikovPetrovich_2012}, and typically describe planetary eccentricity evolution only in the limit $e\to 0$, except for \citet{PapaloizouLarwood_2000}.

Fourth, an alternative and perhaps somewhat under-appreciated approach is to forgo any simplifications of the linearized and Fourier-decomposed equations altogether and to pursue their direct numerical solution in 2D, as first performed by \cite{KorycanskyPollack_1993} for planets on circular orbits. This calculation has also been extended to 3D in \citet{TanakaEtAl_2002}, \citet{TanakaWard_2004}, \citet{TanakaOkada_2024}. In the linear limit this semi-analytical approach should give an {\it exact} solution for the planetary torque \citep{RafikovPetrovich_2012,Petrovich2012} and the Fourier mode structure, allowing one to reconstruct the full spatial distribution of the planet-driven perturbation \citep{MirandaRafikov_2019a,MirandaRafikov_2020}.

In this paper we follow \citet{KorycanskyPollack_1993} and employ this latter, semi-analytical approach, making no assumptions or simplifications apart from the linearity of planet-disc coupling, to explore the orbital evolution of planets in a 2D disc over a range of planetary eccentricity, not necessarily assuming $e$ to be very small. In particular, this allows us to probe a problematic regime of transonic planetary motion for which there is some controversy in the existing results \citep{PapaloizouLarwood_2000,CresswellNelson_2006,IdaEtAl_2020}; namely when $e$ is of the order of the disc aspect ratio. We explore the transonic transition over a whole grid of disc parameters --- aspect ratio, surface density and temperature slopes --- and provide accurate prescriptions for the planetary migration and eccentricity evolution timescales as a function of these parameters and eccentricity. 

Our calculations utilize the semi-analytic framework for studying the eccentric planet-disc coupling that was developed in \citet[][hereafter \paper1]{FairbairnRafikov_2022} to explore the spatial structure of the perturbation induced in a 2D disc by an eccentric planet. This framework was shown to provide a very good match for the time-resolved morphology of the wake induced in the disc by an eccentric planet that was previously investigated via direct hydrodynamic simulations by \citet{ZhuZhang_2022}. Whilst in \paper1 we focused on synthesising the wake morphology, we now exploit this framework to investigate the planet-disc torque acting on an eccentric planet and to provide prescriptions for the evolution of the planetary orbit. 

This paper is organized as follows. In Section \ref{sec:linear_method} we will summarise the prerequisite linear theory and numerical pipeline developed in \paper1. In Section \ref{sec:circular_benchmark} we benchmark this procedure on circular orbits before turning our attention to eccentric orbits in Section \ref{sec:eccentric_torques}. We will discuss these results in the context of previous studies in Section \ref{sec:discussion} before concluding in Section \ref{sec:conclusion}. Furthermore, Appendix \ref{app:parameter_space} describes how to access and process our data which can be used to easily obtain planetary migration and eccentricity evolution timescales over the range of parameters explored in this work. 














